\uuid{bdVH}
\titre{Lecture des paramètres d'une loi normale}
\auteur{Maxime Nguyen}
\organisation{AMSCC}
\datecreate{2026-11-24}
\niveau{L2}
\module{Probabilité}
\chapitre{Probabilité continue}
\sousChapitre{Loi normale}
\difficulte{1}
\theme{loi normale, écart-type, variance, standardisation}

\contenu{
\begin{question}
On note $\mathcal{N}(\mu,\sigma)$ la loi normale d'espérance $\mu$ et d'écart-type $\sigma$, et $\mathbb{V}(X)$ la variance de $X$. Soit $X \sim \mathcal{N}(5, 2)$. Quelles affirmations sont vraies ?
\begin{answers}
\bad{$\mathbb{V}(X) = 2$.}
\good{$\mathbb{V}(X) = 4$.}
\good{$\dfrac{X - 5}{2}$ suit la loi $\mathcal{N}(0,1)$.}
\bad{$\dfrac{X - 5}{4}$ suit la loi $\mathcal{N}(0,1)$.}
\bad{$\dfrac{X - 2}{5}$ suit la loi $\mathcal{N}(0,1)$.}
\end{answers}
\begin{explanations}
Le second paramètre est l'écart-type :
$\sigma = 2$ et $\mathbb{V}(X) = \sigma^2 = 4$. On centre puis on divise par
l'écart-type.

Contre-sens visés : confondre écart-type et variance (A, D) ; intervertir
les rôles des deux paramètres (E). Attention : d'autres ouvrages notent
$\mathcal{N}(\mu, \sigma^2)$ ; il faut toujours vérifier la convention.
\end{explanations}
\end{question}
}
